\section{Background and Workload Characterization}
\label{sec:background}

\subsection{Depth Completion and the Encoder--Decoder Pipeline}

Depth completion estimates a dense metric-depth map from an RGB image and an
incomplete depth observation.  The input depth may contain holes, outliers, and
material-dependent noise, whereas the output must retain both absolute scale
and local geometry for downstream robotic perception.  Modern methods commonly
pair a representation encoder with a dense prediction decoder.  The encoder
extracts global visual and geometric features; the decoder fuses and upsamples
these features to the target depth resolution.

\system{} preserves this algorithmic decomposition and uses its boundary as a
system partition.  The selected depth-completion network and its trained
weights are the workload being accelerated, while this paper focuses on the
low-precision deployment, heterogeneous execution, and dedicated decoder
hardware.  The final manuscript will state the exact backbone, decoder,
input/output resolutions, and feature tensor exported at the partition.

\subsection{Computational Characteristics}

The two stages expose different opportunities for edge acceleration.  The
encoder contains Transformer blocks whose matrix operations and MLP layers can
use Thor's mature GPU and Tensor Core execution stack.  The decoder is composed
primarily of structured dense-prediction operators for which customized
parallelism, buffering, and precision can be implemented on FPGA.  This
difference motivates a stage-level mapping rather than repeatedly moving
individual operators between devices.

The argument must be grounded in measurements.  We therefore profile the
unmodified model on Thor and report the latency of preprocessing, encoder,
decoder, and postprocessing under a fixed resolution, precision, batch size,
and power mode.  We additionally characterize the dominant operator classes,
their arithmetic intensity or compute/memory demand, and the size of the
encoder--decoder boundary tensor.

\subsection{Edge SoC--FPGA Heterogeneous Execution}

An edge SoC offers a complete software environment and high-throughput tensor
units, but its compute and energy budget is shared with sensing, planning, and
control.  An FPGA provides configurable precision, dataflow, and memory
organization for a stable subgraph.  The combination is useful only when the
partition is coarse enough to limit communication and when execution across
devices is scheduled to hide, rather than merely add, transfer latency.

Prior GPU depth-completion systems establish the accuracy of large models, and
prior FPGA accelerators demonstrate the efficiency of specialized low-precision
datapaths.  \system{} connects these directions by retaining the encoder on the
SoC, specializing the decoder on FPGA, and treating the model boundary and
cross-frame schedule as first-class parts of the design.
